\subsection{Estimating a pose from point correspondences (Kabsch / Umeyama)}\label{sec:rot:estimating-pose-point}
\begin{deriv}{closed-form least-squares rigid alignment}
Given model points $p_i$ and measured points $q_i\approx Rp_i+t$, minimise $\sum_i\|q_i-Rp_i-t\|^2$. Setting the derivative w.r.t.\ $t$ to zero: $t=\bar q-R\bar p$ (centroids must align). With centred points $p_i'=p_i-\bar p$, $q_i'=q_i-\bar q$, the cost becomes $\text{const}-2\sum_i(q_i')\T Rp_i'=\text{const}-2\tr(RH)$ with $H=\sum_ip_i'(q_i')\T$. Take the SVD $H=USV\T$; then $\tr(RH)=\tr\big((V\T RU)S\big)$ and $V\T RU$ is orthogonal, so the trace is largest when $V\T RU=I$:
\[
R=V\,\mathrm{diag}\big(1,1,\det(VU\T)\big)\,U\T,\qquad t=\bar q-R\bar p.
\]
The $\det$ term prevents returning a reflection when the data are noisy or nearly planar.
\end{deriv}
\begin{example}{object pose from five 3-D points}
A 20\,cm object has 5 known model points; a depth camera measures them with 2\,mm noise. Kabsch recovers the pose with rotation error $1.30^\circ$ and translation error $2.2$\,mm (RMS residual 3.8\,mm). Note the large rotation error: angular accuracy is roughly noise$/$size, so for small objects orientation is the weak part of the estimate. The same routine is the inner step of ICP, aligns estimated and ground-truth trajectories before computing ATE, and (with a scale factor, Umeyama) aligns monocular SLAM trajectories.
\end{example}
